\documentclass[UTF8,a4paper]{ctexbook}
\usepackage{amsmath,amssymb,amsthm}
\usepackage{graphicx,booktabs,needspace,tikz}
\usepackage[hidelinks]{hyperref}
\hypersetup{pdftitle={从全等变换到谱定理：七门数学的例题讲义},bookmarksnumbered=true}
\setcounter{tocdepth}{1}

\newcounter{example}[chapter]
\renewcommand{\theexample}{\thechapter.\arabic{example}}
\newenvironment{example}[1]{%
  \par\Needspace{5\baselineskip}\medskip
  \refstepcounter{example}\noindent\textbf{例 \theexample\quad #1}\par\nobreak
  \noindent\ignorespaces}{\par}
\newenvironment{solution}{%
  \par\smallskip\noindent\textbf{解析：}\ignorespaces}{\par\medskip}
\newtheorem{theorem}{定理}[chapter]

\title{从全等变换到谱定理\\[.6em]\Large 七门数学的例题讲义}
\author{}
\date{}

\begin{document}
\frontmatter
\maketitle
\tableofcontents
\chapter{导读：一个矩阵，两件事}

这本书从平面上的两条椭圆出发，依次用数学分析、高等代数、近世代数、泛函分析、实分析、复分析、范畴论理解谱定理。导读和高等代数的基础例题从坐标与向量出发；数学分析章会用到基本微积分，后续各章先介绍做题所需的概念，再进入抽象推导。每道例题都给出解析。第一次阅读，可以先看本导读和高等代数章，再按兴趣回到其他各章。七章各有六道例题，另有两道导读题和四道综合题，共四十八道。

\section{先认清要证明哪两个图形全等}

给定实对称矩阵 $A$，记 $q_A(x)=x^TAx$。若一个正交矩阵 $Q$ 满足
\[
Q^TAQ=D=\operatorname{diag}(\lambda_1,\ldots,\lambda_n),\qquad Q^TQ=I,
\]
则 $x=Qu$ 同时给出两条等式：
\begin{align*}
q_A(Qu)&=u^TQ^TAQu=u^TDu,\\
\|Qu-Qv\|^2&=(u-v)^TQ^TQ(u-v)=\|u-v\|^2.
\end{align*}
第一条等式给出等值集的对应：$u^TDu=c$ 当且仅当 $q_A(Qu)=c$。第二条等式说明任意两点的距离在对应中不变。由于 $Q^{-1}=Q^T$，对应是双射，故两个等值图形全等。这一结论也适用于不等式 $q_A(x)\le c$ 所定义的点集。

实对称矩阵的谱定理保证这样的 $Q$ 存在。这里，$A$ 决定二次型，$Q$ 才负责搬动图形。若把两种坐标都画在同一个固定坐标系里，$Q$ 表示主动移动；若只把坐标轴换成 $Q$ 的列，几何点没有移动，改变的是同一点的坐标。

\begin{example}{主例：把标准椭圆转成斜椭圆}
令
\[
A_0=\begin{pmatrix}5&2\\2&2\end{pmatrix},\qquad
Q_0=\frac1{\sqrt5}\begin{pmatrix}2&-1\\1&2\end{pmatrix},\qquad
D_0=\begin{pmatrix}6&0\\0&1\end{pmatrix}.
\]
验证 $Q_0$ 给出 $6u^2+v^2=1$ 与 $5x^2+4xy+2y^2=1$ 的全等对应，并找出对应的一对长轴端点。
\end{example}
\begin{solution}
$Q_0$ 的两列长度都是 $1$，内积是 $(-2+2)/5=0$，故 $Q_0^TQ_0=I$。分别用 $A_0$ 乘两列，得到原列的 $6$ 倍和 $1$ 倍，因此 $A_0Q_0=Q_0D_0$。于是 $Q_0^TA_0Q_0=D_0$。取
\[
x=\frac{2u-v}{\sqrt5},\qquad y=\frac{u+2v}{\sqrt5}.
\]
这就是两椭圆之间的保距双射。长轴端点 $(0,\pm1)$ 映到 $\pm(-1,2)/\sqrt5$，变换前后的端点距离都是 $2$。短轴端点 $(\pm1/\sqrt6,0)$ 映到 $\pm(2,1)/\sqrt{30}$，短轴长度均为 $2/\sqrt6$。而 $\det Q_0=1$，所以这里的保距变换是旋转，转角为 $\arctan(1/2)$。
\end{solution}

\begin{figure}[htbp]
\centering
\begin{tikzpicture}[x=1.55cm,y=1.55cm,>=stealth]
\begin{scope}
\draw[->] (-1.22,0)--(1.3,0) node[right] {$u$};
\draw[->] (0,-1.2)--(0,1.22) node[above] {$v$};
\draw[thick,domain=0:360,samples=101,smooth,variable=\t]
  plot ({cos(\t)/sqrt(6)},{sin(\t)});
\draw[dashed] (0,1)--({1/sqrt(6)},0)--(0,-1)--cycle;
\fill (0,1) circle (1.1pt) node[above right] {$R$};
\fill (0,-1) circle (1.1pt) node[below right] {$S$};
\fill ({1/sqrt(6)},0) circle (1.1pt) node[above right] {$P$};
\node at (0,-1.46) {$6u^2+v^2=1$};
\end{scope}
\draw[->] (1.7,0.5)--(2.55,0.5) node[midway,above] {$Q_0$};
\begin{scope}[xshift=6.5cm]
\draw[->] (-1.22,0)--(1.3,0) node[right] {$x$};
\draw[->] (0,-1.2)--(0,1.22) node[above] {$y$};
\draw[thick,domain=0:360,samples=101,smooth,variable=\t]
  plot ({(2*cos(\t)/sqrt(6)-sin(\t))/sqrt(5)},
        {(cos(\t)/sqrt(6)+2*sin(\t))/sqrt(5)});
\coordinate (rp) at ({-1/sqrt(5)},{2/sqrt(5)});
\coordinate (sp) at ({1/sqrt(5)},{-2/sqrt(5)});
\coordinate (pp) at ({2/sqrt(30)},{1/sqrt(30)});
\draw[dashed] (rp)--(pp)--(sp)--cycle;
\fill (rp) circle (1.1pt) node[above left] {$R'$};
\fill (sp) circle (1.1pt) node[below right] {$S'$};
\fill (pp) circle (1.1pt) node[above right] {$P'$};
\node at (0,-1.46) {$5x^2+4xy+2y^2=1$};
\end{scope}
\end{tikzpicture}
\caption{两图比例相同。椭圆和内接三角形由同一个正交矩阵逐点对应。}
\end{figure}

\section{什么条件不能省}

\begin{example}{把椭圆拉成圆，为什么不能说全等}
令 $s=\sqrt6\,u,\ t=v$，则 $6u^2+v^2=1$ 变成 $s^2+t^2=1$。说明这个代换为什么不能用于证明全等。
\end{example}
\begin{solution}
这个代换把短轴端点间的距离 $2/\sqrt6$ 拉长到 $2$。它的矩阵为 $P=\operatorname{diag}(\sqrt6,1)$，而
\[
P^TP=\operatorname{diag}(6,1)\ne I.
\]
如果有其他平面等距变换把圆送到椭圆，圆心的像必到椭圆上的每个点等距。椭圆的两对对称端点迫使这样的候选点只能是原点，但原点到长轴端点与短轴端点的距离分别为 $1$ 与 $1/\sqrt6$。所以圆和此椭圆不全等。一般可逆代换可以化简方程，几何全等要求保距。
\end{solution}

\section{符号与阅读路线}

$\langle x,y\rangle$ 表示内积，$\|x\|$ 表示范数；有限维实空间中转置用 $T$，复空间共轭转置用 $*$。符号 $\sigma(T)$ 在无限维章节表示算子的谱：使 $\lambda I-T$ 没有有界逆的所有复数 $\lambda$。它可能包含并非特征值的点。

\begin{center}
\begin{tabular}{lp{.62\textwidth}}
\toprule
视角 & 要回答的问题\\
\midrule
数学分析 & 哪个方向最陡？旋转前后的曲线长度和曲率怎样比较？\\
高等代数 & 如何找正交特征基？特征值能判定哪一类全等？\\
近世代数 & 群作用如何把矩阵分成轨道？哪些变换不改变它？\\
泛函分析 & 无穷维里何时还有正交特征基？如何逐项计算？\\
实分析 & 旋转怎样改变积分？保测与保距有什么区别？\\
复分析 & 旋转、共形和全纯有什么关系？围道积分能读出哪些谱信息？\\
范畴论 & 换坐标以后，哪些构造与同构相容？\\
\bottomrule
\end{tabular}
\end{center}

\mainmatter
\chapter{数学分析：从变化率读出主轴}
\label{ch:calculus}

同一个二次型，沿不同方向走，增长速度不同。分析的办法是限制方向、求导、找极值，再观察旋转究竟保留了哪些量。本章沿用导读中的 $A_0,Q_0,D_0$。

\section{在单位圆上寻找特征方向}

设单位向量 $e(\theta)=(\cos\theta,\sin\theta)^T$。把二次型限制在单位圆上，得到单变量函数 $r(\theta)=q_A(e(\theta))$。若 $A$ 对称，求导得 $r'(\theta)=2e'(\theta)^TAe(\theta)$。由于 $e,e'$ 是互相垂直的两个单位向量，$r'(\theta)=0$ 恰好说明 $Ae$ 沿 $e$ 方向，即 $e$ 是特征向量。

\begin{example}{用三角恒等式找最大值和最小值}
在 $x^2+y^2=1$ 上，求 $5x^2+4xy+2y^2$ 的最大值、最小值与相应方向。
\end{example}
\begin{solution}
代入 $x=\cos\theta,\ y=\sin\theta$，再用倍角公式，得到
\[
r(\theta)=\frac72+\frac32\cos2\theta+2\sin2\theta
         =\frac72+\frac52\cos(2\theta-\phi),
\quad \cos\phi=\frac35,\quad\sin\phi=\frac45.
\]
故最大值为 $6$，最小值为 $1$。取 $\alpha=\arctan(1/2)$，则 $\cos2\alpha=3/5,\ \sin2\alpha=4/5$，可取 $\phi=2\alpha$。最大方向为 $\theta=\alpha$ 或 $\alpha+\pi$，对应 $\pm(2,1)^T/\sqrt5$；最小方向与它垂直，为 $\pm(-1,2)^T/\sqrt5$。在方向 $e$ 上，点 $te$ 满足 $q_A(te)=t^2q_A(e)$。当等值常数为 $1$ 时，端点到原点的距离是 $1/\sqrt{q_A(e)}$，所以增长最快的方向反而是短轴方向。
\end{solution}

\begin{example}{拉格朗日乘子为什么会变成特征值}
在 $x^2+y^2+z^2=1$ 上求 $4x^2+4xy+4y^2+z^2$ 的极值，并解释乘子的意义。
\end{example}
\begin{solution}
约束 $g(x,y,z)=x^2+y^2+z^2=1$ 上 $\nabla g=2(x,y,z)^T\ne0$。极值点上目标函数沿球面切方向的导数为零，所以 $\nabla f$ 与 $\nabla g$ 平行。写 $\nabla f=\lambda\nabla g$，约去 $2$ 后得
\[
(4-\lambda)x+2y=0,\quad 2x+(4-\lambda)y=0,\quad(1-\lambda)z=0.
\]
前两式相加、相减，得 $(6-\lambda)(x+y)=0$ 与 $(2-\lambda)(x-y)=0$。由单位长度知 $x+y,x-y,z$ 至少一个非零，故候选 $\lambda$ 为 $6,2,1$。分别取单位向量 $(1,1,0)/\sqrt2$、$(1,-1,0)/\sqrt2$、$(0,0,1)$ 都能达到。因为 $Ae=\lambda e$ 时 $e^TAe=\lambda\|e\|^2=\lambda$，乘子恰是相应函数值。球面紧致，极值必达到，故最大值为 $6$，最小值为 $1$。
\end{solution}

\section{二阶导数怎样辨认临界点}

设二元函数 $f$ 在临界点 $a$ 的邻域二阶连续可微，则其二阶 Taylor 展开为
\[
f(a+h)=f(a)+\tfrac12h^THh+o(\|h\|^2),
\qquad H=(\partial_i\partial_j f(a)).
\]
二阶偏导连续时，Hessian 矩阵 $H$ 对称。把它正交对角化，便能看清局部的增减方向。余项 $o(\|h\|^2)$ 除以 $\|h\|^2$ 后趋于零，因此在各特征值严格同号时，它不会改变充分靠近临界点的符号判断。

\begin{example}{带有交叉项和三次项的严格极小值}
判断 $f(x,y)=5x^2+4xy+2y^2+x^3$ 在原点的局部性质，并给出一个明确的邻域。
\end{example}
\begin{solution}
原点是临界点，$H=2A_0$ 的特征值为 $12,2$。令 $(x,y)^T=Q_0(u,v)^T$，则二次部分为 $6u^2+v^2\ge u^2+v^2=x^2+y^2$。设 $r=\sqrt{x^2+y^2}$，有 $|x|^3\le r^3$，所以
\[
f(x,y)-f(0,0)\ge r^2-r^3=r^2(1-r)>0\qquad (0<r<1).
\]
这就直接证明了严格局部极小值。三次项一般会改变整条等值曲线，以上论证用的是二阶近似及余项控制，不能宣称 $f=1$ 与二次型椭圆全等。
\end{solution}

\begin{example}{同一个退化 Hessian 对应三种情况}
比较 $f_+(x,y)=x^2+y^4$、$f_-(x,y)=x^2-y^4$ 和 $f_0(x,y)=x^2$ 在原点的性质。
\end{example}
\begin{solution}
三个函数在原点的梯度都为零，Hessian 都是 $\operatorname{diag}(2,0)$，谱完全相同。$f_+$ 的两项非负，只有原点同时为零，故是严格极小；$f_-$ 在 $x$ 轴上为正、在 $y$ 轴上为负，故是鞍点；$f_0$ 在整条 $y$ 轴上为零，故是非严格极小。特征值 $0$ 对应的方向没有二阶信息，必须检查更高阶项或直接比较函数值。
\end{solution}

\section{旋转保留曲线的长度与弯曲程度}

对分段一阶连续可微的参数曲线 $\gamma(t)$，切向量是 $\gamma'(t)$，弧长是速度的积分。对二次可微且速度不为零的平面曲线，曲率大小为
\[
\kappa(t)=\frac{|\det(\gamma'(t),\gamma''(t))|}{\|\gamma'(t)\|^3}.
\]
在正交变换下，速度范数不变，行列式只可能改变符号，故弧长与曲率大小都不变。

\begin{example}{不必算出椭圆周长，也能比较周长}
令 $\gamma(t)=(a\cos t,b\sin t)$，$a=1/\sqrt6,b=1$。证明标准椭圆与 $Q_0$ 旋转后的椭圆周长相等，并比较对应的四分之一段弧长。
\end{example}
\begin{solution}
$\gamma'(t)=(-a\sin t,b\cos t)$，故
\[
\|\gamma'(t)\|=\sqrt{a^2\sin^2t+b^2\cos^2t}.
\]
旋转后的参数曲线是 $\eta(t)=Q_0\gamma(t)$，其导数为 $Q_0\gamma'(t)$，速度大小相同。于是任意参数区间 $[s,t]$ 上的弧长相等，特别是在 $[0,\pi/2]$ 上，从短轴端点到相邻长轴端点的弧长都等于
\[
\int_0^{\pi/2}\sqrt{\tfrac16\sin^2t+\cos^2t}\,dt.
\]
两条完整曲线都绕一周，积分区间 $[0,2\pi]$ 的弧长相同；无须把积分求成初等函数。
\end{solution}

\begin{example}{哪一个轴端点更弯}
计算上述标准椭圆在 $(a,0)$ 和 $(0,b)$ 的曲率大小，再给出斜椭圆中对应点的曲率。
\end{example}
\begin{solution}
有 $\gamma''(t)=(-a\cos t,-b\sin t)$，因此 $\det(\gamma',\gamma'')=ab$。所以
\[
\kappa(t)=\frac{ab}{(a^2\sin^2t+b^2\cos^2t)^{3/2}}.
\]
在 $t=0$ 处为 $a/b^2=1/\sqrt6$；在 $t=\pi/2$ 处为 $b/a^2=6$。长轴端点 $(0,1)$ 更弯。斜椭圆对应点依次为 $(2,1)/\sqrt{30}$ 和 $(-1,2)/\sqrt5$。因为 $\det Q_0=1$ 且速度不变，曲率仍为 $1/\sqrt6$ 和 $6$。若用镜面反射代替旋转，带符号曲率会反号，但曲率大小保持不变。
\end{solution}

\chapter{高等代数：用正交特征基证明全等}
\label{ch:linear}

对称性是谱定理的关键。若 $A^T=A$，则 $\langle Ax,y\rangle=\langle x,Ay\rangle$。若 $Av=\lambda v,\ Aw=\mu w$ 且 $\lambda\ne\mu$，这个等式给出 $(\lambda-\mu)\langle v,w\rangle=0$，即两个特征方向垂直。但这还没有保证特征向量够用，下面给出完整论证。

\begin{theorem}[实对称矩阵的谱定理]
每个实对称矩阵 $A$ 都有一组标准正交的特征基。把它们作为列组成 $Q$，便有 $Q^TAQ=D$，且 $Q^TQ=I$。
\end{theorem}
\begin{proof}
单位球面紧致，连续函数 $x^TAx$ 在其上取得最大值。设单位向量 $v$ 是一个最大点。任取单位向量 $w\perp v$，函数
\[
t\longmapsto(\cos t\,v+\sin t\,w)^TA(\cos t\,v+\sin t\,w)
\]
在 $t=0$ 的导数是 $2w^TAv$，必须为零。因此 $Av$ 垂直于 $v^\perp$，即 $Av=\lambda v$。对任意 $z\perp v$，有 $v^TAz=(Av)^Tz=\lambda v^Tz=0$，故 $A$ 把 $v^\perp$ 映回自身，且限制在其中仍是对称的。在这个维数少一的空间重复论证；一维时结论直接成立。这就得到所需的正交特征基。
\end{proof}

\section{从系数算出主轴和距离}

\begin{example}{从原方程开始算出旋转矩阵}
不借用导读的答案，计算 $5x^2+4xy+2y^2$ 的正交标准形，并验证两对对应点之间的距离。
\end{example}
\begin{solution}
交叉项的系数 $4$ 分给矩阵的两个非对角位置，各为 $2$。于是
\[
\det(A_0-\lambda I)=(5-\lambda)(2-\lambda)-4
=(\lambda-6)(\lambda-1).
\]
当 $\lambda=6$ 时，$-x+2y=0$，可取方向 $(2,1)$；当 $\lambda=1$ 时，$2x+y=0$，可取 $(-1,2)$。两者的长度均为 $\sqrt5$，除以长度得到 $Q_0$ 的两列。由于 $A_0Q_0=Q_0D_0$ 且 $Q_0^TQ_0=I$，得到标准形 $6u^2+v^2$。

例如在标准椭圆取 $P=(1/\sqrt6,0)$ 和 $R=(0,1)$，距离平方是 $1/6+1=7/6$。它们的像是 $P'=(2,1)/\sqrt{30}$、$R'=(-1,2)/\sqrt5$；把坐标差平方相加，交叉项相消，仍为 $7/6$。这个计算核验了一对点；两椭圆全等的证明需要对任意 $p,r$ 使用 $\|Q_0(p-r)\|=\|p-r\|$，再用 $Q_0^{-1}$ 证明一一对应。
\end{solution}

\begin{example}{负特征值不妨碍等值图形全等}
把 $4xy=1$ 化成无交叉项的方程，并说明换元后得到的是什么图形。
\end{example}
\begin{solution}
系数矩阵 $B=\begin{pmatrix}0&2\\2&0\end{pmatrix}$。向量 $(1,1)^T$ 与 $(-1,1)^T$ 的特征值分别为 $2,-2$。令
\[
Q=\frac1{\sqrt2}\begin{pmatrix}1&-1\\1&1\end{pmatrix},
\qquad
x=\frac{u-v}{\sqrt2},\quad y=\frac{u+v}{\sqrt2}.
\]
代入得 $4xy=2u^2-2v^2$。因此原方程与 $2u^2-2v^2=1$ 的双曲线全等。后者的顶点是 $(\pm1/\sqrt2,0)$，送回原坐标为 $\pm(1/2,1/2)$；渐近线 $u=\pm v$ 对应原来的两条坐标轴。正负号决定图形的类型；能否全等对应由 $Q$ 是否正交决定。
\end{solution}

\section{重特征值与投影}

\begin{example}{在二维特征空间中找两条正交轴}
设 $w=(1,1,1)^T$，$B=I+ww^T$。求正交标准形，并找出 $x^TBx=1$ 的三个半轴长度。
\end{example}
\begin{solution}
由于 $w^Tw=3$，有 $Bw=w+w(w^Tw)=4w$。若 $w^Tx=0$，则 $Bx=x$，所以垂直于 $w$ 的平面是特征值 $1$ 的特征空间。取
\[
v_1=\frac{(1,1,1)^T}{\sqrt3},\quad
v_2=\frac{(1,-1,0)^T}{\sqrt2},\quad
v_3=\frac{(1,1,-2)^T}{\sqrt6}.
\]
逐项内积可验证正交且都是单位向量。取这些列得到 $Q$，于是 $Q^TBQ=\operatorname{diag}(4,1,1)$。对应椭球方程 $4u_1^2+u_2^2+u_3^2=1$，三个半轴长为 $1/2,1,1$。后两根轴可在同一平面内一起转动；不变的是该平面与半径，而不是某两根指定的轴。
\end{solution}

\begin{example}{不选特征向量的正负号，也能投影}
只用 $A_0$ 的矩阵与特征值 $6,1$，求投影到特征值 $6$ 的正交投影矩阵。将 $e=(1,0)^T$ 分成两个正交分量，计算 $A_0^2e$。
\end{example}
\begin{solution}
多项式 $p(t)=(t-1)/5$ 在 $6$ 处取 $1$、在 $1$ 处取 $0$。因此 $p(A_0)$ 在两个特征方向上分别乘 $1,0$，故
\[
P_6=\frac{A_0-I}{5}=\frac15\begin{pmatrix}4&2\\2&1\end{pmatrix},
\qquad
P_1=I-P_6=\frac15\begin{pmatrix}1&-2\\-2&4\end{pmatrix}.
\]
$P_6e=(4,2)^T/5$，$P_1e=(1,-2)^T/5$，两者内积为零，和为 $e$。作用两次等于把这两个分量分别乘 $36,1$，得到
\[
A_0^2e=36P_6e+P_1e=(29,14)^T.
\]
这与直接求 $A_0(5,2)^T$ 一致。投影矩阵不随单位特征向量的正负号改变；后面将用复积分和范畴论再得到同一投影。
\end{solution}

\section{什么时候可以用谱判定全等}

\begin{example}{反过来：全等的正定单位椭圆一定同谱吗}
设 $A,B$ 都是实对称正定的二阶矩阵。证明 $x^TAx=1$ 和 $y^TBy=1$ 全等，当且仅当 $A,B$ 的特征值（计重数）相同。
\end{example}
\begin{solution}
若两矩阵同谱，把各自的单位特征基按相同顺序排列，写成 $A=Q_ADQ_A^T$、$B=Q_BDQ_B^T$。则 $U=Q_BQ_A^T$ 正交，且 $U^TBU=A$，所以 $y=Ux$ 给出保距对应。

反过来，任一平面刚体等距变换写成 $y=Ux+b$，其中 $U$ 正交。椭圆的中心唯一：若一个非空有界点集有两个中心，则两次中心对称的复合是非零平移；反复平移会产生无界点列，矛盾。两椭圆中心都是原点，故 $b=0$。任取 $v\ne0$，令 $x=v/\sqrt{v^TAv}$，则 $x$ 在第一条椭圆上，故 $(Ux)^TB(Ux)=1$。乘回分母得到 $v^TAv=v^TU^TBUv$。这对所有 $v$ 成立；分别令 $v=e_i$、$e_i+e_j$，可读出两边对称矩阵的各项系数相等，所以 $A=U^TBU$。正交相似保特征值。
\end{solution}

\begin{example}{同谱的非对称矩阵能否照搬这个结论}
比较 $C=I_2$ 和 $J=\begin{pmatrix}1&1\\0&1\end{pmatrix}$。它们是否同谱、相似？对应的实二次型单位等值集是否全等？
\end{example}
\begin{solution}
两者的特征多项式都是 $(\lambda-1)^2$，但 $J-I$ 的核只有一维，$C-I$ 的核是二维。相似保持核的维数，因此它们不相似，当然也不正交相似。另一个直接理由是任意可逆 $S$ 都有 $S^{-1}IS=I$。

对实向量 $x$，$x^TJx=x^T\bigl((J+J^T)/2\bigr)x$，因为反对称部分在两侧乘同一向量时为零。$J$ 的对称部分为 $\begin{pmatrix}1&1/2\\1/2&1\end{pmatrix}$，特征值是 $3/2,1/2$。它的单位等值椭圆半轴长为 $\sqrt{2/3}$ 和 $\sqrt2$；$C$ 的单位等值集是半径 $1$ 的圆。由上一题可知不全等。对称性不能省；一般实矩阵的二次型由对称部分决定。
\end{solution}

\chapter{近世代数：正交群怎样作用在二次型上}
\label{chap:abstract}

把两个图形比较一次，只需找到一个保距变换。若要研究\emph{所有}可能的转动、反射和摆放方式，就需要能把这些变换组织起来的语言。群描述变换的复合；群作用描述这些变换怎样改变对象。本章只考虑固定原点的变换，把二次型的系数矩阵作为对象，看同一个中心椭圆可以有多少种摆法，以及哪些变换使它保持原状。

\section{旋转与反射的复合}
一个\textbf{群}具有封闭的乘法，乘法满足结合律，有单位元，且每个元素都有逆元。所有满足 \(U^TU=I\) 的 \(n\) 阶实矩阵在矩阵乘法下组成正交群 \(O(n)\)。因为
\[
 (UV)^T(UV)=V^TU^TUV=I,\qquad U^{-1}=U^T,
\]
两个正交变换可以连续做，也可以各自撤销。由 \((\det U)^2=1\) 可知 \(\det U=\pm1\)。其中行列式为 \(1\) 的矩阵构成特殊正交群 \(SO(n)\)。

在平面上，设 \(R_\theta\) 是逆时针旋转 \(\theta\) 的矩阵，设
\[
 R_\theta=\begin{pmatrix}\cos\theta&-\sin\theta\\
 \sin\theta&\cos\theta\end{pmatrix},\qquad
 S=\begin{pmatrix}1&0\\0&-1\end{pmatrix}.
\]
若 \(U\in O(2)\) 的第一列为 \((\cos\theta,\sin\theta)^T\)，第二列只能是 \((-\sin\theta,\cos\theta)^T\) 或它的相反向量。故 \(U\) 或为 \(R_\theta\)，或为 \(R_\theta S\)。前者是旋转；后者是先关于横轴反射，再旋转。

\begin{example}{两次反射，为什么转过两倍角}
先关于横轴反射，再关于与横轴成 \(30^\circ\) 的直线反射。求合成变换，并求 \((2,0)^T\) 的像。若交换反射次序，结果怎样？
\end{example}
\begin{solution}
先约定角度的正方向：\(\beta\) 从正横轴逆时针量起，表示\textbf{反射直线的方向角}，不是点被旋转的角。直线没有指定朝向，所以 \(\beta\) 与 \(\beta+180^\circ\) 表示同一条反射线；反射矩阵也应相同。以下固定坐标轴、主动移动点，并使用列向量。\textbf{矩阵乘积中，右边的变换先做。}

\textbf{第一步：写出关于斜线的反射。}先把点旋转 \(-\beta\)，使反射直线转到横轴；关于横轴反射；最后旋转 \(\beta\) 回到原来的位置。因此
\[
 F_\beta=R_\beta S R_{-\beta}
 =\begin{pmatrix}
 \cos 2\beta&\sin 2\beta\\
 \sin 2\beta&-\cos 2\beta
 \end{pmatrix}.
\]
矩阵中的 \(2\beta\) 并不表示 \(F_\beta\) 是旋转：它的行列式为 \(-1\)，而且整条反射直线上的点都不动。若一个非零向量原方向角为 \(\varphi\)，反射后方向角为 \(2\beta-\varphi\)。

\textbf{第二步：按题目次序复合。}横轴反射把方向角 \(\varphi\) 变为 \(-\varphi\)，再作斜线反射得到 \(2\beta+\varphi\)，所以每个方向都统一增加 \(2\beta\)。矩阵计算也给出
\[
 F_\beta S
 =\begin{pmatrix}
 \cos 2\beta&-\sin 2\beta\\
 \sin 2\beta&\cos 2\beta
 \end{pmatrix}
 =R_{2\beta}.
\]
交换次序则为
\[
 SF_\beta
 =\begin{pmatrix}
 \cos 2\beta&\sin 2\beta\\
 -\sin 2\beta&\cos 2\beta
 \end{pmatrix}
 =R_{-2\beta}.
\]
这两个合成变换互为逆变换。因为每次反射都做两遍就回到变换前的位置，即 \(S^2=F_\beta^2=I\)，所以撤销 \(F_\beta S\) 必须按相反次序做 \(SF_\beta\)。

\textbf{第三步：代入 \(\beta=30^\circ\)。}此时
\[
 F_{30^\circ}=\frac12\begin{pmatrix}1&\sqrt3\\\sqrt3&-1\end{pmatrix},
 \quad F_{30^\circ}S=R_{60^\circ},
 \quad SF_{30^\circ}=R_{-60^\circ}.
\]
对题目中的具体点，两种次序的每一步分别为
\begin{align*}
 (2,0)^T&\xrightarrow{\ S\ }(2,0)^T
          \xrightarrow{\ F_{30^\circ}\ }(1,\sqrt3)^T,\\
 (2,0)^T&\xrightarrow{\ F_{30^\circ}\ }(1,\sqrt3)^T
          \xrightarrow{\ S\ }(1,-\sqrt3)^T.
\end{align*}
两像的长度都是 \(\sqrt{1+3}=2\)，但点的位置不同。反向合成把正向的像还原：\(R_{-60^\circ}(1,\sqrt3)^T=(2,0)^T\)。

\textbf{第四步：核对保距。}直接相乘可得 \(F_\beta^TF_\beta=I\) 及 \(R_\theta^TR_\theta=I\)。所以对任意两点 \(p,q\)，各次反射及两个合成变换 \(T\) 都满足
\[
 \|T(p)-T(q)\|^2
 =(p-q)^TT^TT(p-q)=\|p-q\|^2.
\]
这是对所有点对成立的保距证明；具体点的长度检查只是其中一个例子。两次反射的合成虽保距，交换次序却通常改变结果。本题的 \(60^\circ\) 与 \(-60^\circ\) 旋转正好展示了这个区别。
\end{solution}

\section{轨道描述所有摆法，稳定子描述对称}
设 \(A\) 是实对称矩阵。定义左作用
\[
 U\cdot A=UAU^T \qquad (U\in O(n)).
\]
它确实满足 \(I\cdot A=A\) 和 \(U\cdot(V\cdot A)=(UV)\cdot A\)。这里的写法对应主动移动图形：若 \(x=Uz\)，则
\[
 x^T(U\cdot A)x=z^TAz.
\]
因此 \(U\) 把 \(A\) 的等值图形旋转或反射到 \(U\cdot A\) 的等值图形。所有 \(U\cdot A\) 组成 \(A\) 的\textbf{轨道}；满足 \(U\cdot A=A\) 的变换组成 \(A\) 的\textbf{稳定子}。谱定理给出每条轨道的对角代表；把特征值按大小排序后，这个对角代表唯一。

\begin{example}{只许转直角，椭圆有几种摆法}
令
\[
 A=\begin{pmatrix}5&2\\2&2\end{pmatrix},\qquad
 J=\begin{pmatrix}0&-1\\1&0\end{pmatrix}.
\]
在只允许 \(I,J,J^2,J^3\) 的情况下，求 \(A\) 的轨道和稳定子，并解释图形的含义。
\end{example}
\begin{solution}
\(J\) 是逆时针旋转 \(90^\circ\)，且 \(J^2=-I\)、\(J^4=I\)。直接计算
\[
 JAJ^T=\begin{pmatrix}2&-2\\-2&5\end{pmatrix}=:B.
\]
由于 \((-I)A(-I)^T=A\)，\(J^2\) 使矩阵回到 \(A\)，而 \(J^3=-J\) 给出 \(B\)。故轨道为 \(\{A,B\}\)，在这四种变换中的稳定子为 \(\{I,-I\}\)。
因此椭圆 \(5x^2+4xy+2y^2=1\) 旋转 \(90^\circ\) 得到不同摆法：例如 \((1/\sqrt5,0)\) 在原椭圆上，代入新方程左边却为 \(2/5\)。旋转 \(180^\circ\) 则回到原状。两个矩阵的迹都为 \(7\)，行列式都为 \(6\)，所以特征值同为 \(6,1\)；这也与正交共轭保谱相符。
\end{solution}

正定二次型的单位等值集足以确定其系数矩阵。事实上，若两个正定二次型 \(q_1,q_2\) 的单位等值集相同，任取 \(x\ne0\)，令 \(z=x/\sqrt{q_1(x)}\)，则 \(q_1(z)=q_2(z)=1\)。把缩放比例代回便得 \(q_1(x)=q_2(x)\)。对实对称矩阵，分别比较各平方项和交叉项的系数，就知道两个矩阵相同。因此，研究下一题的矩阵等式，正好等于研究相应椭圆的摆放。

\begin{example}{求出把标准椭圆移到斜椭圆的所有正交矩阵}
设 \(D=\operatorname{diag}(6,1)\)，并令
\[
 Q=\frac1{\sqrt5}\begin{pmatrix}2&-1\\1&2\end{pmatrix},
 \qquad A=QDQ^T.
\]
求所有满足 \(UDU^T=A\) 的 \(U\in O(2)\)。其中有几个是旋转？
\end{example}
\begin{solution}
先求 \(D\) 的稳定子。令 \(V=Q^TU\)，则 \(UDU^T=A\) 等价于 \(VDV^T=D\)，即 \(VD=DV\)。写 \(V=(v_{ij})\)，比较非对角元得
\[
 v_{12}=6v_{12},\qquad 6v_{21}=v_{21}.
\]
故 \(V\) 是对角矩阵。正交性又要求 \(v_{11},v_{22}\in\{1,-1\}\)。反过来，每个这样的 \(V\) 都固定 \(D\)，所以全部解为
\[
 U=Q\operatorname{diag}(\varepsilon_1,\varepsilon_2),
 \qquad \varepsilon_1,\varepsilon_2\in\{1,-1\}.
\]
这四个矩阵对应两个主轴方向可以各自反向。因为 \(\det Q=1\)，其中行列式为 \(1\) 的恰是 \(Q\) 和 \(-Q\)，两种旋转相差 \(180^\circ\)。其余两个是反射。椭圆的主轴方向由两条特征直线确定，因此改变特征向量的正负号不会改变主轴。
\end{solution}

\section{特征值重合以后，对称变多了}
\begin{example}{同一特征值的两个方向可以混合吗}
令 \(D_3=\operatorname{diag}(6,1,1)\)。求所有满足 \(UD_3U^T=D_3\) 的 \(U\in O(3)\)，并说明 \(6x^2+y^2+z^2=1\) 的相应对称。
\end{example}
\begin{solution}
由 \(UD_3=D_3U\) 得，对于第 \(i,j\) 个矩阵元 \(u_{ij}\)，有 \((d_j-d_i)u_{ij}=0\)，其中 \(d_1=6,d_2=d_3=1\)。所以第一行、第一列的非对角元必须为零，而右下角的 \(2\times2\) 块可以有非零的非对角元。再用正交性得到
\[
 U=\begin{pmatrix}\varepsilon&0\\0&W\end{pmatrix},
 \qquad \varepsilon=\pm1,\quad W\in O(2).
\]
每个这样的 \(U\) 都与 \(D_3\) 交换，故这就是全部稳定子。几何上，\(x\) 轴可以反向，\(yz\) 平面内可以任意旋转或反射。比如固定 \(x\)，把 \((y,z)\) 换成 \((-z,y)\)，表达式 \(y^2+z^2\) 保持不变。与上一题相比，特征值 \(1\) 的特征空间由一条直线变成一个平面，允许在其中任取一组正交单位基。
\end{solution}

\section{区分代数合同和欧氏全等}
一般可逆换元 \(x=Pz\) 给出 \(A\mapsto P^TAP\)，称为合同。当 \(P\) 正交时，\(P^T=P^{-1}\)，这个公式也可读作相似变换。一般合同保留二次型的惯性，即正、负、零平方项的个数；正交共轭还保留特征值的具体数值。这个差别决定了变换是否保留图形的轴长。

\begin{example}{惯性相同的两条双曲线一定全等吗}
比较 $2x^2-8y^2=1$ 与 $u^2-v^2=1$。找出一个合同换元，再判断两条双曲线是否全等。
\end{example}
\begin{solution}
取 $P=\operatorname{diag}(1/\sqrt2,1/\sqrt8)$，即 $x=u/\sqrt2,\ y=v/\sqrt8$。代入得到 $u^2-v^2=1$；两边都是一个正平方项和一个负平方项，惯性相同。但 $P$ 不保距离。要证明不存在其他等距变换，可比较两个连通分支之间的最小距离。第一条曲线的右支满足 $x\ge1/\sqrt2$，左支满足 $x\le-1/\sqrt2$。两支之间任意两点的距离至少为 $\sqrt2$，在 $(\pm1/\sqrt2,0)$ 处达到。第二条曲线的同一距离是 $2$，在 $(\pm1,0)$ 处达到。等距双射保持连通分支及分支间的最小距离，因此两条双曲线不全等。此例表明，惯性确定二次型能化成哪一类符号标准形，却不记录几何尺度。
\end{solution}

另一方面，仅凭某一个等值点集相同，未必能推出对应矩阵在同一正交轨道。例如 \(x^2-2y^2=0\) 与 \(2x^2-4y^2=0\) 都给出直线 \(x=\pm\sqrt2\,y\)，但对应特征值分别为 \(\{1,-2\}\) 和 \(\{2,-4\}\)。在零水平上整体乘非零常数不改变点集，却可能改变谱。

\section{有限群的平均：给定对称，反过来找内积}
若有限群 \(G\) 的每个元素 \(g\) 对应一个可逆矩阵 \(\rho(g)\)，且 \(\rho(gh)=\rho(g)\rho(h)\)、\(\rho(e)=I\)，便得到一个矩阵表示。原来的欧氏内积未必被它保留。把所有矩阵作用后的内积取平均，可得到新内积 \(x^THy\)，其中
\[
 H=\frac{1}{|G|}\sum_{g\in G}\rho(g)^T\rho(g).
\]
对非零 \(x\)，和式中的单位元项已给出正贡献，所以 \(H\) 正定。对任意 \(h\in G\)，右乘 \(h\) 会重新排列 \(G\) 中的元素，故
\[
 \rho(h)^T H\rho(h)
 =\frac1{|G|}\sum_{g\in G}\rho(gh)^T\rho(gh)=H.
\]
谱定理允许取 \(C=H^{1/2}\)，于是 \(C^TC=H\)。用坐标 \(z=Cx\) 表示时，各 \(\rho(h)\) 变成正交矩阵 \(C\rho(h)C^{-1}\)。

\begin{example}{把一个斜反射写成普通反射}
令 \(T=\begin{pmatrix}1&2\\0&-1\end{pmatrix}\)。验证 \(\{I,T\}\) 是群，求它平均得到的 \(H\)，并找坐标 \(z=Cx\)，使 \(T\) 在新坐标下是关于横轴的反射。
\end{example}
\begin{solution}
矩阵相乘得 \(T^2=I\)，且 \(T\ne I\)，故有一个二元群。原欧氏度量下，\(T(0,1)^T=(2,-1)^T\)，长度从 \(1\) 变为 \(\sqrt5\)，所以 \(T\) 不正交。平均得到
\[
 H=\frac{I+T^TT}{2}
  =\begin{pmatrix}1&1\\1&3\end{pmatrix}.
\]
取 \(C=\begin{pmatrix}1&1\\0&\sqrt2\end{pmatrix}\)，便有 \(C^TC=H\)。对 \(x=(a,b)^T\)，新坐标是 \(z=(a+b,\sqrt2 b)^T\)；而
\[
 CTx=C(a+2b,-b)^T=(a+b,-\sqrt2 b)^T.
\]
所以新坐标下恰为 \((z_1,z_2)\mapsto(z_1,-z_2)\)。在原坐标中，\(T\) 保持的新长度是 \(\sqrt{x^THx}=\sqrt{(a+b)^2+2b^2}\)。这说明一个有限群的表示可以在合适的内积下保距；改变内积与证明它原本就保持欧氏距离是两件事。
\end{solution}

\chapter{泛函分析：函数空间里的保距与对角化}
\label{chap:functional}

把平面中的两个坐标推广成无穷多个坐标，就得到平方可和序列空间
\begin{align*}
 \ell^2&=\left\{x=(x_1,x_2,\ldots):\sum_{n=1}^{\infty}|x_n|^2<\infty\right\},\\
 \|x\|^2&=\sum_{n=1}^{\infty}|x_n|^2.
\end{align*}
也可以用函数当向量，例如在 \(L^2[0,1]\) 中用 \(\|f\|^2=\int_0^1|f(t)|^2\,dt\) 衡量长度。两个空间都是希尔伯特空间：有内积，并且按内积给出的范数是完备的。本章内积的计算用实数；谈算子的谱时，可以在相应的复空间中作同样的计算。

设 \(T\) 是有界线性算子，若 \(\langle Tx,y\rangle=\langle x,Ty\rangle\) 对所有 \(x,y\) 成立，就称 \(T\) 自伴。若 \(U\) 是满射线性等距，则 \(U^{-1}=U^*\)；在实空间中相当于正交变换，在复空间中称酉变换。令 \(B=U^*TU\)，则
\[
 \langle Bz,z\rangle=\langle TUz,Uz\rangle,
 \qquad \|Uz-Uw\|=\|z-w\|.
\]
因此，\(U\) 把 \(B\) 的二次型等值集保距地映到 \(T\) 的相应等值集。这一推导与维数无关。难点是：能否找到由特征向量组成的正交基？

\section{从两个坐标推广到无穷多个坐标}

\begin{example}{只旋转前两个坐标}
在 \(\ell^2\) 中，令 \(T\) 在前两个坐标上按矩阵 \(A=\begin{pmatrix}5&2\\2&2\end{pmatrix}\) 作用，其余坐标变成零。令 \(U\) 在前两个坐标上按 \(Q=\frac1{\sqrt5}\begin{pmatrix}2&-1\\1&2\end{pmatrix}\) 作用，其余坐标照旧。求 \(U^*TU\)，说明两算子的二次型等值集如何对应。
\end{example}
\begin{solution}
把 \(z\) 写成 \((z_1,z_2,z_3,\ldots)\)。由于 \(Q^TQ=I\)，前两个坐标的平方和保持不变，后面的平方和也不变，所以 \(U\) 保范数。用 \(Q^T\) 恢复前两个坐标，其余照旧，给出 \(U\) 的逆。因此 \(U\) 是满射等距。矩阵计算 \(Q^TAQ=\operatorname{diag}(6,1)\) 只涉及前两个坐标，故
\[
 U^*TU(z_1,z_2,z_3,\ldots)=(6z_1,z_2,0,\ldots).
\]
在新坐标中，等值 \(1\) 的条件是 \(6z_1^2+z_2^2=1\)，其余 \(z_n\) 只需平方可和。前两项描出椭圆，后面的坐标可自由变化，得到一个无限维的柱形点集。用 \(U\) 映射，就得到原坐标中 \(5x_1^2+4x_1x_2+2x_2^2=1\) 的点集。自由坐标没有被丢掉；两个点集之间保距且一一对应。
\end{solution}

\begin{example}{傅里叶投影怎样去掉第二谐波}
在实函数空间 \(L^2[0,2\pi]\) 中取内积
\[
\langle f,g\rangle=\frac{1}{2\pi}\int_0^{2\pi}f(t)g(t)\,dt.
\]
令 \(P\) 为到 \(\operatorname{span}\{1,\cos t,\sin t\}\) 的正交投影。对 \(f(t)=3+2\cos t+\sin(2t)\)，求 \(Pf\)、误差平方 \(\|f-Pf\|^2\)，并指出 \(P\) 的特征子空间。
\end{example}
\begin{solution}
由三角函数乘积化和差公式，不同整数频率的正弦、余弦彼此正交，且非零频率的正弦、余弦积分为零。归一化后，
\[
 e_0=1,\qquad e_1=\sqrt2\cos t,\qquad e_2=\sqrt2\sin t
\]
是一组单位正交向量。因此投影是把 \(f\) 的三个系数 \(\langle f,e_j\rangle\) 乘回相应的 \(e_j\) 再相加。这三个系数为 \(3,\sqrt2,0\)，于是
\begin{align*}
 Pf&=3+2\cos t, & f-Pf&=\sin(2t),\\
 \|Pf\|^2&=9+2=11, & \|f-Pf\|^2&=\frac12.
\end{align*}
在目标子空间上，\(P\) 等于恒等算子，特征值是 \(1\)；在其正交补上，\(P\) 等于零，特征值是 \(0\)。于是 \(P\) 既自伴又满足 \(P^2=P\)。这里的「对角化」就是把任意 \(f\) 分成 \(Pf\) 与 \(f-Pf\) 两个互相垂直的分量，勾股关系给出 \(\|f\|^2=11+1/2\)。投影本身一般会缩短向量，不是等距变换。
\end{solution}

\section{紧自伴算子：什么时候有特征向量基}

紧算子把有界序列送到具有收敛子列的序列。有限秩有界算子是紧的，紧算子关于算子范数取极限后仍然紧。紧自伴谱定理保证：在可分希尔伯特空间中，可把非零特征值的正交特征向量和核空间的正交基合起来，得到整个空间的正交基；若有无穷多个非零特征值（计重数），它们只能趋向零。下面直接算两个例子。

\begin{example}{趋于零的对角元素}
在 \(\ell^2\) 中定义 \((Tx)_n=x_n/n\)。证明 \(T\) 自伴且紧，求它的谱。解释为什么零在谱中，却不是特征值。
\end{example}
\begin{solution}
逐项相乘可得 \(\|Tx\|\leq\|x\|\)，并且
\(\langle Tx,y\rangle=\sum_n x_n\overline{y_n}/n=\langle x,Ty\rangle\)。
令 \(T_N\) 只保留前 \(N\) 个输出坐标。它有有限秩，且
\[
 \|(T-T_N)x\|^2
 =\sum_{n>N}\frac{|x_n|^2}{n^2}
 \leq \frac{\|x\|^2}{(N+1)^2}.
\]
取 \(x=e_{N+1}\) 达到等号，所以 \(\|T-T_N\|=1/(N+1)\to0\)，证明 \(T\) 紧。

第 \(n\) 个标准基向量 \(e_n\) 满足 \(Te_n=e_n/n\)。若 \(\lambda\) 不在闭集 \(\{0,1,1/2,\ldots\}\) 中，则 \(\inf_n|1/n-\lambda|>0\)。逐坐标除以 \(1/n-\lambda\) 就是 \(T-\lambda I\) 的有界逆。反过来，\(\lambda=1/n\) 时有非零核；\(\lambda=0\) 时，若存在有界逆，则
\[
 1=\|e_n\|=\|T^{-1}Te_n\|\leq \|T^{-1}\|/n,
\]
对大 \(n\) 矛盾。因此谱恰为 \(\{0,1,1/2,\ldots\}\)。但 \(Tx=0\) 强迫每个 \(x_n=0\)，所以零不对应任何非零特征向量。谱记录的是有界可逆性，范围比特征值更广。
\end{solution}

\begin{example}{把积分方程变成二元方程}
在 \(L^2[0,1]\) 中设 \(u(s)=1\)、\(v(s)=\sqrt3(2s-1)\)，定义
\begin{align*}
 K(s,t)&=5u(s)u(t)+2u(s)v(t)+2v(s)u(t)+2v(s)v(t),\\
 (Tf)(s)&=\int_0^1K(s,t)f(t)\,dt.
\end{align*}
求 \(T\) 的非零特征值、对应的单位特征函数，并解积分方程 \(f+Tf=u+v\)。
\end{example}
\begin{solution}
直接积分，得到 \(\int_0^1 u^2=1\)、\(\int_0^1uv=0\)、\(\int_0^1 v^2=1\)。若 \(a=\langle f,u\rangle\)、\(b=\langle f,v\rangle\)，把积分拆成四项可得
\[
 Tf=(5a+2b)u+(2a+2b)v.
\]
因此 \(T\) 在 \(\operatorname{span}\{u,v\}\) 上的矩阵正是 \(A\)，在正交补上为零。两组方向互相垂直，且 \(A=A^T\)，故 \(T\) 自伴；值域至多二维，故 \(T\) 紧。由同一个 \(Q\) 可得
\[
 \phi_6=(2u+v)/\sqrt5,\qquad
 \phi_1=(-u+2v)/\sqrt5,\qquad
 T\phi_6=6\phi_6,\quad T\phi_1=\phi_1.
\]
积分方程的右端在二维子空间里，对等式取正交补分量，可知 \(f\) 的该分量为零。又
\[
 u+v=\frac3{\sqrt5}\phi_6+\frac1{\sqrt5}\phi_1.
\]
于是解方程只需把两个系数分别除以 \(1+6\) 和 \(1+1\)：
\[
 f=\frac3{7\sqrt5}\phi_6+\frac1{2\sqrt5}\phi_1
 =\frac1{14}u+\frac27v.
\]
代回得 \(Tf=\frac{13}{14}u+\frac57v\)，所以 \(f+Tf=u+v\)。在正交补上 \(I+T=I\)，在两个特征方向上乘数均不为零，故解唯一。复杂的积分形式之所以能化简，是因为算子的作用集中在一个二维平面内。
\end{solution}

\section{无穷维中需要重新检查的条件}

\begin{example}{保距离的右移为什么不是换基}
令 \(S(x_1,x_2,\ldots)=(0,x_1,x_2,\ldots)\)。证明 \(S\) 保距离，求 \(S^*\)，并说明 \(S^*S=I\) 为何不能推出 \(SS^*=I\)。对前例的对角算子 \((Tx)_n=x_n/n\)，再算 \(S^*TS\)。
\end{example}
\begin{solution}
插入一个零不改变平方和，所以 \(\|Sx-Sy\|^2=\|x-y\|^2\)。由
\[
 \langle Sx,y\rangle
 =\sum_{n\geq1}x_n\overline{y_{n+1}}
 =\langle x,(y_2,y_3,\ldots)\rangle
\]
得 \(S^*(y_1,y_2,\ldots)=(y_2,y_3,\ldots)\)。先插零后删第一项，得到 \(S^*S=I\)；先删第一项再插零，得到
\[
 SS^*(y_1,y_2,\ldots)=(0,y_2,y_3,\ldots).
\]
这说明值域是第一坐标为零的闭子空间，\(e_1\) 没有原像。计算得
\[
TSx=(0,x_1/2,x_2/3,\ldots).
\]
再用 \(S^*\) 去掉第一个坐标，得到
\[
 S^*TSx=(x_1/2,x_2/3,\ldots).
\]
其对角元素从 \(1/2\) 开始，原来属于 \(T\) 的特征值 \(1\) 消失了。此处的保距只发生在 \(\ell^2\) 到该子空间之间；对 \(T\) 的二次型而言，这是限制到该子空间再换坐标，不能当作整个空间上的正交相似。
\end{solution}

\begin{example}{每个特征值都为正，等值集仍可能无界}
对 \((Tx)_n=x_n/n\)，设 \(E=\{x:\langle Tx,x\rangle=1\}\)。证明 \(E\) 无界。再取 \(\varepsilon>0\)，求 \(E_\varepsilon=\{x:\langle (T+\varepsilon I)x,x\rangle=1\}\) 上 \(\|x\|\) 的最小值与上确界。
\end{example}
\begin{solution}
有 \(\langle Tx,x\rangle=\sum_n |x_n|^2/n\)。对非零 \(x\) 该式总为正，但没有统一的正常数 \(m\) 使它对所有 \(x\) 都不小于 \(m\|x\|^2\)。事实上，\(x=\sqrt n\,e_n\) 属于 \(E\)，而其范数 \(\sqrt n\) 可任意大。

对 \(x\in E_\varepsilon\)，由 \(0<\langle Tx,x\rangle\leq\|x\|^2\) 可得
\[
 \varepsilon\|x\|^2<1\leq (1+\varepsilon)\|x\|^2,
 \qquad \frac1{\sqrt{1+\varepsilon}}\leq\|x\|
 <\frac1{\sqrt{\varepsilon}}.
\]
取 \(x=e_1/\sqrt{1+\varepsilon}\) 达到下界；取 \(x_n=e_n/\sqrt{\varepsilon+1/n}\)，范数趋近右端，所以右端是上确界但取不到。有限维中，所有特征值为正便有正的最小特征值；无穷维中，即使算子在每个非零方向都给出正的二次型值，也要另查是否有统一的严格正下界。
\end{solution}

本章三个现象应分开：正交换坐标使二次型等值集保距对应；紧自伴条件提供特征向量基；谱的统一正下界控制等值集的大小。一般有界自伴算子的谱分解还能用投影来表达，下一章给出没有特征向量的例子。

\chapter{实分析：从全等图形到测度与连续谱}
\label{chap:real-analysis}

实分析可以把「图形有多大」推广为勒贝格测度，也可以在 \(L^2\) 中研究函数在各点的取值。测度与距离是两种不同的信息。若 \(Q^TQ=I\)，则 \((\det Q)^2=1\)，由线性变量代换公式，对可测集 \(E\subseteq\mathbb R^n\) 有
\[
 m(QE)=\int_E |\det Q|\,dx=m(E).
\]
平移同样保测。因而欧氏全等推出相同的测度；面积相等本身不足以推出全等。

\section{保面积与保距离}

\begin{example}{相同的面积能否推出全等}
取矩形 \(R=[0,2]\times[0,1]\)。分别作 \(Q=\frac1{\sqrt5}\begin{pmatrix}2&-1\\1&2\end{pmatrix}\) 与剪切 \(S=\begin{pmatrix}1&2\\0&1\end{pmatrix}\)。比较 \(QR\)、\(SR\) 与 \(R\) 的面积、全等关系。
\end{example}
\begin{solution}
两矩阵的行列式都等于 \(1\)，因此变量代换给出
\[
 m(QR)=m(SR)=m(R)=2.
\]
由于 \(Q^TQ=I\)，它使任意两点的距离保持不变，又可逆，所以 \(QR\) 与 \(R\) 全等。剪切矩阵把 \(R\) 的两个顶点 \((0,0),(2,1)\) 分别送到 \((0,0),(4,1)\)。在 \(R\) 中，两点横坐标之差的绝对值至多 \(2\)，纵坐标之差的绝对值至多 \(1\)，因此任意两点距离都不超过 \(\sqrt5\)。然而 \(SR\) 里已有相距 \(\sqrt{17}\) 的两点，故不存在把 \(R\) 保距地映到整个 \(SR\) 的映射。这不仅说明 \(S\) 本身不保距，还排除了用别的等距变换使两个图形全等的可能。
\end{solution}

\begin{example}{椭圆区域和两条等值线之间的面积}
令 \(q_A(x,y)=5x^2+4xy+2y^2\)，\(E_c=\{(x,y):q_A(x,y)\leq c\}\)。求 \(c>0\) 时的面积，以及 \(0<a<b\) 时区域 \(\{a<q_A\leq b\}\) 的面积。说明等值线 \(q_A=c\) 的平面测度为零。
\end{example}
\begin{solution}
作 \((x,y)^T=Q(u,v)^T\)，则 \(q_A=6u^2+v^2\)，且 \(Q\) 不改变面积。再令
\[
 u=\sqrt{c/6}\,s,\qquad v=\sqrt c\,t.
\]
此时 \(6u^2+v^2\leq c\) 等价于 \(s^2+t^2\leq1\)。第二次变换的行列式为 \(c/\sqrt6\)，所以
\[
 m(E_c)=\frac{\pi c}{\sqrt6},\qquad
 m(\{a<q_A\leq b\})=m(E_b)-m(E_a)=\frac{\pi(b-a)}{\sqrt6}.
\]
这里 \(E_a\subseteq E_b\)，故可直接相减。对 \(0<h<c\)，等值线 \(q_A=c\) 包含在 \(\{c-h<q_A\leq c\}\) 中，其面积至多 \(\pi h/\sqrt6\)；让 \(h\) 趋零即得结论。第一步的 \(Q\) 是全等变换；第二步是用来算面积的伸缩，不保持全等。
\end{solution}

\section{把积分与函数的取值分布联系起来}

\begin{example}{利用等值区域计算二维高斯积分}
求 \(\int_{\mathbb R^2}e^{-q_A(x,y)}\,dx\,dy\)。若 \(c>0\)，在 \(E_c\) 中按面积均匀取一点 \(X\)，求 \(q_A(X)/c\) 的分布。
\end{example}
\begin{solution}
对每个点，\(q_A\geq0\)，而 \(e^{-q_A}=\int_{q_A}^\infty e^{-s}\,ds\)。把积分下限改成 \(0\) 并乘上条件 \(q_A\leq s\) 的示性函数，再用非负函数的 Tonelli 定理交换积分次序：
\[
 \int_{\mathbb R^2}e^{-q_A}\,dx\,dy
 =\int_0^\infty e^{-s}m(E_s)\,ds
 =\frac{\pi}{\sqrt6}\int_0^\infty se^{-s}\,ds
 =\frac{\pi}{\sqrt6}.
\]
最后一步由分部积分得到 \(\int_0^\infty se^{-s}\,ds=1\)。均匀取点时，对 \(0\leq r\leq1\)，有
\[
 \mathbb P(q_A(X)/c\leq r)=\frac{m(E_{rc})}{m(E_c)}=r.
\]
\(r=0\) 时只有原点，面积为零；\(r=1\) 时整个区域都满足条件。因此 \(q_A(X)/c\) 在 \([0,1]\) 上均匀分布。计算依靠的是正交变换下的面积不变和面积随 \(c\) 的线性增长，而非点的角度均匀。
\end{solution}

\begin{example}{旋转函数图样，什么数据随之保持}
定义 \(Uf(z)=f(Qz)\)，其中 \(f\in L^2(\mathbb R^2)\)。证明 \(U\) 是满射等距。对 \(f=2\mathbf 1_{E_1}-\mathbf 1_{E_2\setminus E_1}\)，求 \(\|Uf\|_2^2\) 和 \(m(\{z:|Uf(z)|>t\})\)（\(t\geq0\)）。相同的取值分布能否推出两个图样旋转或平移后重合？
\end{example}
\begin{solution}
变量代换 \(x=Qz\) 给出
\[
 \|Uf\|_2^2=\int_{\mathbb R^2}|f(Qz)|^2\,dz
 =\int_{\mathbb R^2}|f(x)|^2\,dx.
\]
若两个函数只在零测集上不同，则拉回后也只在零测集上不同，故 \(U\) 在 \(L^2\) 的等价类上定义良好。\(U^{-1}g(x)=g(Q^Tx)\) 给出满射性。并且
\[
 \{z:|Uf(z)|>t\}=Q^T\{x:|f(x)|>t\};
\]
两集的测度相等。由于 \(m(E_1)=m(E_2\setminus E_1)=\pi/\sqrt6\)，可得
\[
 \|Uf\|_2^2=\frac{5\pi}{\sqrt6},\qquad
 m(\{|Uf|>t\})=
 \begin{cases}
 2\pi/\sqrt6,&0\leq t<1,\\
 \pi/\sqrt6,&1\leq t<2,\\
 0,&t\geq2.
 \end{cases}
\]
反推不成立。例如 \(R=[0,2]\times[0,1]\) 与
\(C=([0,1]\cup[3,4])\times[0,1]\) 面积相同，示性函数取值都只有 \(0,1\)，相应正阈值的超水平集测度相同。但 \(R\) 的直径是 \(\sqrt5\)，\(C\) 的直径是 \(\sqrt{17}\)；等距变换保直径，所以无法使两个图样重合。即使允许相差零测集也不行：任一顶点的任意邻域内都有正面积的矩形部分，删掉零测集仍可任意逼近该顶点，因此两集的直径各自保持不变。
\end{solution}

\section{乘法算子：谱可能是一整个区间}

在 \(L^2[0,1]\) 中定义 \(Mf(t)=tf(t)\)。它像对角矩阵一样「每个位置单独乘一个数」，但在函数空间里，单点的测度为零，不能把每个位置当成一个非零坐标向量。此时谱分解应看成按取值范围切分函数。

\begin{example}{自伴算子没有非零特征向量}
证明 \(M\) 有界自伴、非紧，且谱为 \([0,1]\)。证明 \(M\) 没有特征值。
\end{example}
\begin{solution}
因为 \(0\leq t\leq1\)，有 \(\|Mf\|_2\leq\|f\|_2\)。内积中乘子 \(t\) 为实数，所以 \(\langle Mf,g\rangle=\langle f,Mg\rangle\)。若 \(Mf=\lambda f\)，则 \((t-\lambda)f(t)=0\) 几乎处处。除至多一个点以外，\(t-\lambda\) 不为零，所以 \(f=0\) 几乎处处；不存在非零特征函数。

若 \(\lambda\notin[0,1]\)，则 \(1/(t-\lambda)\) 在 \([0,1]\) 上有界，乘以它就得到 \(M-\lambda I\) 的有界逆。若 \(\lambda\in[0,1]\)，取 \(I_h=[0,1]\cap(\lambda-h,\lambda+h)\)（\(0<h<1\)），令 \(f_h=\mathbf1_{I_h}/\sqrt{m(I_h)}\)。则
\[
 \|f_h\|_2=1,\qquad \|(M-\lambda I)f_h\|_2\leq h.
\]
若有有界逆，便会推出 \(1\leq\|(M-\lambda I)^{-1}\|h\)，让 \(h\) 趋零即矛盾。因此每个 \(\lambda\in[0,1]\) 都属于谱。

再取 \([1/2,1]\) 中两两不交、长度为正的区间 \(J_n=(1/2+2^{-(n+2)},1/2+2^{-(n+1)})\)，其中 \(n\geq1\)。令 \(g_n=\mathbf1_{J_n}/\sqrt{m(J_n)}\)。各 \(g_n\) 的范数为 \(1\)，各 \(Mg_n\) 的范数至少为 \(1/2\)，且支撑两两不交。因此 \(n\ne k\) 时
\[
 \|Mg_n-Mg_k\|_2^2=\|Mg_n\|_2^2+\|Mg_k\|_2^2\geq\frac12.
\]
序列 \((Mg_n)\) 无收敛子列，故 \(M\) 非紧。此例说明无穷维有界自伴算子不必能按普通特征向量逐个对角化。
\end{solution}

\begin{example}{把一个谱区间上的分量取出来}
对上述 \(M\)，给定可测集 \(B\subseteq[0,1]\)，定义 \(P_B f=\mathbf1_B f\)。验证它是正交投影。令 \(f(t)=\sqrt3\,t\)，求 \(\|P_{[a,b]}f\|_2^2\) 与 \(\langle Mf,f\rangle\)。再证明 \(M\) 与 \(M_{1-t}:f(t)\mapsto(1-t)f(t)\) 的二次型等值集可由满射等距对应。
\end{example}
\begin{solution}
由 \(\mathbf1_B^2=\mathbf1_B\) 可知 \(P_B^2=P_B\)，并且 \(f=P_Bf+(f-P_Bf)\)，两项支撑互不重叠而正交。因此 \(P_B\) 正是到「在 \(B\) 外为零」的闭子空间的正交投影。对 \(0\leq a\leq b\leq1\)，
\begin{align*}
 \|P_{[a,b]}f\|_2^2&=\int_a^b 3t^2\,dt=b^3-a^3,\\
 \|f\|_2^2&=1,\qquad
 \langle Mf,f\rangle=\int_0^1 3t^3\,dt=\frac34.
\end{align*}
可把 \(3t^2\,dt\) 看作此单位向量在各段谱区间上的「平方长度分配」：单点分量为零，整个区间却可以有正的分量。这是连续谱中用投影表达分解的具体方式。

令 \(Vf(t)=f(1-t)\)。换元 \(s=1-t\) 给出 \(\|Vf\|_2=\|f\|_2\)，且 \(V^{-1}=V\)。逐点计算得
\[
 (V^{-1}MVf)(t)=(1-t)f(t),
 \qquad
 \langle M_{1-t}f,f\rangle=\langle MVf,Vf\rangle.
\]
所以 \(V\) 把前一二次型的任一等值集保距地送到后一二次型的相应等值集，反向用 \(V^{-1}\) 即可。整个过程不需要选出一个特征向量。全等关系来自满射等距所给出的共轭；是否有普通特征向量基是另一个问题。
\end{solution}

\chapter{复分析：保角、保距与复积分选谱}
\label{chap:complex}

把实平面上的点 $(x,y)$ 记成复数 $z=x+iy$，两点距离就是 $|z-w|$。于是，旋转与反射都可以用复数公式计算。复分析还提供了另一种看谱定理的方式：把矩阵的特征值当作复平面上的点，用一条闭曲线圈出想要的谱成分。读本章时先把“是否保距”理解成检验 $|f(z)-f(w)|=|z-w|$，再看复可微性和复积分各自能说明什么。

\section{用复变函数辨别平面变换}

复导数 $f'(z)$ 由差商 $[f(z+h)-f(z)]/h$ 在复数 $h$ 趋于零时的极限给出。与实数情形不同，$h$ 可沿平面中的任意方向趋于零，极限必须相同。一个函数在某开集的每一点都复可微，就称它在该开集上全纯。

\begin{example}{用复数写出椭圆主轴的旋转}
设 $Q=\frac1{\sqrt5}\begin{pmatrix}2&-1\\1&2\end{pmatrix}$。把 $Q$ 写成复平面上的映射，求 $1+i$ 的像，并证明这个映射保距。
\end{example}
\begin{solution}
对 $z=x+iy$，直接相乘：
\[
 \frac{2+i}{\sqrt5}z=\frac{2x-y}{\sqrt5}+i\frac{x+2y}{\sqrt5}.
\]
实部、虚部恰好是 $Q(x,y)^T$ 的两个分量。系数 $a=(2+i)/\sqrt5$ 的模为 $1$，辐角为 $\arctan(1/2)$，故这是逆时针旋转。$1+i$ 的像为 $(1+3i)/\sqrt5$。任取两点 $z,w$，
\[
 |az-aw|=|a|\,|z-w|=|z-w|.
\]
因此一对点的距离和整条曲线的形状大小都保持。再加一个固定复数 $b$ 也保距，因为计算差时 $b$ 抵消。
\end{solution}

\begin{example}{关于直线的反射为何不是全纯函数}
把点关于直线 $y=x$ 作反射。写出复数形式，计算 $2+i$ 的像，并判断它在何处复可微。
\end{example}
\begin{solution}
反射交换两个坐标，因此 $R(x+iy)=y+ix=i\overline{x+iy}$，$R(2+i)=1+2i$。对任意两点，
\[
 |R(z)-R(w)|=|i\overline{z-w}|=|z-w|.
\]
但在任一点 $z$，差商为 $i\bar h/h$。取非零实数 $h$ 趋于 $0$，差商恒为 $i$；取纯虚数 $h=it$ 趋于 $0$，差商恒为 $-i$。两个方向的极限不同，所以处处不复可微。由于其共轭函数 $\overline{R(z)}=-iz$ 全纯，$R$ 称为反全纯映射。一般而言，关于过原点、与实轴夹角为 $\theta$ 的直线反射，可以先旋转 $-\theta$、再共轭、最后旋转 $\theta$，所得公式为 $e^{2i\theta}\bar z$。
\end{solution}

\begin{example}{局部保角为何不推出保距}
考察 $f(z)=z^2$ 在 $z_0=1+i$ 的作用。比较相互垂直的切向方向 $1$、$i$ 的像，并说明此映射能否保持任意两点距离。
\end{example}
\begin{solution}
$f(z_0+h)-f(z_0)=2z_0h+h^2$，故一阶近似为乘以 $2+2i$。两个切向量分别变成 $2+2i$ 与 $-2+2i$；对应实向量的内积为 $2(-2)+2\cdot2=0$。它们都逆时针旋转 $\pi/4$，长度乘 $2\sqrt2$。一般地，全纯函数在 $f'(z_0)\ne0$ 处，切向量都乘同一个非零复数 $f'(z_0)$，因而保持相交曲线的有向夹角。这是局部共形性。相比之下，$1,2$ 的距离为 $1$，它们的像 $1,4$ 的距离为 $3$，故 $f$ 不是欧氏等距映射。在 $z=0$，$f'(0)=0$，也不能套用上述保角论证。
\end{solution}

\begin{example}{把单位圆盘映回自身，也可能改变距离}
设 $a=\frac12$，$g(z)=\frac{z-a}{1-az}$。证明 $g$ 把开单位圆盘 $\{z:|z|<1\}$ 一一映到自身，且全纯、保角。检验它是否保持欧氏距离。
\end{example}
\begin{solution}
分母在圆盘内不为零。因 $a$ 是实数，两个模的平方分别为 $1-a(z+\bar z)+a^2|z|^2$ 与 $|z|^2-a(z+\bar z)+a^2$，相减得
\[
 |1-az|^2-|z-a|^2=(1-a^2)(1-|z|^2)>0.
\]
因此 $|g(z)|<1$。从 $w=g(z)$ 解得 $z=(w+a)/(1+aw)$；把 $a$ 换成 $-a$，同一个恒等式证明逆映射也把圆盘映入自身，所以 $g$ 是双射。它的导数
\[
 g'(z)=\frac{1-a^2}{(1-az)^2}
\]
处处非零，故在圆盘内保角。可是 $g(0)=-1/2$、$g(-1/2)=-4/5$，对应距离由 $1/2$ 变成 $3/10$。集合映到自身、局部保角与保持任意两点距离，是三件需要分别验证的事。
\end{solution}

\section{用柯西积分选出特征方向}

在以下两例中，$A=\begin{pmatrix}5&2\\2&2\end{pmatrix}$ 作用在 $\mathbb C^2$ 上。把实矩阵看成复矩阵之后，$Q^*=Q^T$（星号表示共轭转置），且 $A=Q\operatorname{diag}(6,1)Q^*$ 仍然成立。对于 $z\notin\{1,6\}$，矩阵 $(zI-A)^{-1}$ 称为预解式。这里 $z$ 是用于圈选特征值的复参数，向量本身则有两个复坐标。

沿逆时针简单闭曲线积分时，矩阵积分按每个元素分别计算。柯西积分公式告诉我们：若 $f$ 在曲线及其内部的一个开邻域上全纯，则 $f(z)/(z-\lambda)$ 的积分除以 $2\pi i$，在 $\lambda$ 位于曲线内部时等于 $f(\lambda)$，在 $\lambda$ 位于外部时等于 $0$。这里假定 $\lambda$ 不在曲线上。把预解式按特征值拆成分式，每个分式就有一个对应的极点，积分通过“圈内有无该极点”把这些分量区分开。

\begin{example}{绕特征值 $6$ 积分，得到秩一投影}
令 $\Gamma$ 为逆时针圆周 $|z-6|=2$。求
\[
 P=\frac{1}{2\pi i}\int_\Gamma(zI-A)^{-1}\,dz,
\]
并判定 $P$ 把向量投影到哪条直线上。
\end{example}
\begin{solution}
直接求逆得到
\[
 (zI-A)^{-1}=\frac{1}{(z-6)(z-1)}
 \begin{pmatrix}z-2&2\\2&z-5\end{pmatrix}.
\]
圆周圈住 $6$ 而不圈住 $1$，矩阵 $\frac1{z-1}\begin{pmatrix}z-2&2\\2&z-5\end{pmatrix}$ 的每个元素在圆盘上全纯。因此柯西积分公式直接把其中的 $z$ 换成 $6$：
\[
 P=\frac15\begin{pmatrix}4&2\\2&1\end{pmatrix}=qq^*,\qquad q=\frac1{\sqrt5}\binom{2}{1}.
\]
$q^*q=1$，故 $P^*=P$、$P^2=P$。$P$ 将 $x$ 变成 $q(q^*x)$，即沿着 $q$ 保留分量，垂直于 $q$ 的分量被抹去。又 $Aq=6q$，所以它是到特征值 $6$ 的特征直线的正交投影，记作 $P_6$。若换成圈住 $1$ 而不圈住 $6$ 的圆周，计算时将分子矩阵除以 $z-6$，再代入 $z=1$，得 $P_1=\frac15\begin{pmatrix}1&-2\\-2&4\end{pmatrix}=I-P_6$。积分所选择的是整个特征子空间，并不依赖于 $q$ 的相位选择。
\end{solution}

\begin{example}{复指数把自伴矩阵变成保距运动}
令 $t\in\mathbb R$，按指数级数定义 $e^{itA}$。用柯西积分写出它，证明它是幺正矩阵，并计算 $t=\pi/5$ 时的结果。
\end{example}
\begin{solution}
由于 $A=6P_6+P_1$，$P_6P_1=P_1P_6=0$，且两者都是幂等矩阵，所以 $A^n=6^nP_6+P_1$ 对所有 $n\ge1$ 成立；$A^0=I=P_6+P_1$。在收敛的指数级数中逐项代入，可得
\[
 e^{itA}=e^{6it}P_6+e^{it}P_1.
\]
同样地，把预解式写成 $P_6/(z-6)+P_1/(z-1)$，对正向圆周 $|z|=7$ 使用柯西公式，得到
\[
 e^{itA}=\frac{1}{2\pi i}\int_{|z|=7}e^{itz}(zI-A)^{-1}\,dz.
\]
由于 $|e^{6it}|=|e^{it}|=1$，两块分别只改变相位，于是 $(e^{itA})^*e^{itA}=P_6+P_1=I$。它保持 $\mathbb C^2$ 中任意两点的距离。取 $t=\pi/5$，$e^{6it}=-e^{it}$，故
\[
 e^{i\pi A/5}=e^{i\pi/5}(I-2P_6)
 =\frac{e^{i\pi/5}}5\begin{pmatrix}-3&-4\\-4&3\end{pmatrix}.
\]
右边的实矩阵因子 $I-2P_6$ 在特征方向 $q$ 上变号，在其垂直方向上保持不动；额外的相位因子也是等距的。原矩阵 $A$ 沿 $q$ 把长度扩大六倍，但经过函数 $\lambda\mapsto e^{it\lambda}$ 后，两块的伸缩模长都变成 $1$。
\end{solution}

\chapter{范畴论：保留结构的箭头与谱投影}
\label{chap:category}

前面一直在做一种比较：两种写法是否描述了同一个几何结构？范畴论把比较所允许的变换称为\textbf{箭头}。先规定对象和允许的箭头，再问哪些对象同构。改变允许的箭头，分类的结论就可能改变。

\section{把正交对角化读成一个同构}

本章使用有限维复内积空间；实数情形只需把“幺正”换成“正交”，把 $U^*$ 换成 $U^T$。称一个线性映射 $U$ 为幺正同构，意思是它是双射且保持内积。由此 $U^{-1}=U^*$，也有 $\|Ux-Uy\|=\|x-y\|$。

定义范畴 $\mathcal H$：对象是内积空间 $H$，箭头是幺正同构。定义范畴 $\mathcal S$：对象是带有自伴算子的内积空间 $(H,A)$，其中 $A^*=A$；从 $(H,A)$ 到 $(K,B)$ 的箭头 $U$ 除了要幺正，还必须满足
\[
 BU=UA.
\]
这个条件称为\textbf{交织}：先作用 $A$ 再换坐标，与先换坐标再作用 $B$，结果相同。恒等映射满足条件。若另有箭头 $V:(K,B)\to(L,C)$，则 $C(VU)=VBU=VUA$，所以复合仍是箭头；$BU=UA$ 也给出 $AU^{-1}=U^{-1}B$，所以逆映射仍是箭头。所有箭头都可逆的范畴叫\textbf{群胚}，因此 $\mathcal H,\mathcal S$ 都是群胚。

\begin{example}{同一个 $Q$ 为什么既保距又交织}
设 $A=\begin{pmatrix}5&2\\2&2\end{pmatrix}$，$D=\operatorname{diag}(6,1)$，$Q=\frac1{\sqrt5}\begin{pmatrix}2&-1\\1&2\end{pmatrix}$。证明 $Q$ 是从 $(\mathbb C^2,D)$ 到 $(\mathbb C^2,A)$ 的箭头，并解释在实平面中的几何意义。
\end{example}
\begin{solution}
$Q$ 的两列都是单位向量，内积为 $(-2+2)/5=0$，因此 $Q^*Q=I$。第一列乘 $A$ 变成自身的 $6$ 倍，第二列乘 $A$ 不变，于是 $AQ=QD$。交换图为
\[
\begin{array}{ccc}
 \mathbb C^2 &\xrightarrow{D}& \mathbb C^2\\
 \big\downarrow\rlap{$\scriptstyle Q$}&&\big\downarrow\rlap{$\scriptstyle Q$}\\
 \mathbb C^2 &\xrightarrow{A}& \mathbb C^2.
\end{array}
\]
两条路径给出的等式是 $AQ=QD$；$Q^*$ 是反向箭头。由于 $Q$ 的元素是实数，它和它的逆都把实向量送到实向量。限制到实平面，令 $x=Qu$，就有 $x^TAx=u^TQ^TAQu=u^TDu$，同时任意点对距离不变。故 $6u_1^2+u_2^2=1$ 的椭圆被保距地送到 $x^TAx=1$ 的椭圆。一般的复幺正矩阵未必保持实平面，例如 $iI$ 会把非零实向量送成纯虚向量。因此，从复空间的箭头回到实几何时，要确认这个实子空间确实被保留。
\end{solution}

\begin{example}{同构分类要记录特征值的重数}
在 $\mathbb C^3$ 中，比较 $D_1=\operatorname{diag}(6,6,1)$、$D_2=\operatorname{diag}(1,6,6)$、$D_3=\operatorname{diag}(6,1,1)$。确定哪些带算子的对象同构，并说明一般判据。
\end{example}
\begin{solution}
设 $e_1,e_2,e_3$ 为标准基。令 $Ue_1=e_2,\ Ue_2=e_3,\ Ue_3=e_1$，则 $U$ 把一组标准正交基送到另一组，故幺正。逐列检查得到 $D_2U=UD_1$，所以前两个对象同构。

若 $BU=UA$，则 $Av=\lambda v$ 推出 $B(Uv)=\lambda Uv$；对 $U^{-1}$ 再作同样的推理，$U$ 就把 $\ker(A-\lambda I)$ 一一映到 $\ker(B-\lambda I)$。$D_1$ 的 $6$ 特征子空间为二维，$D_3$ 的为一维，故两者不同构，尽管它们不计重数的特征值集合都是 $\{1,6\}$。

反之，若每个 $\lambda$ 的特征子空间维数相同，分别在对应的子空间中取标准正交基，并把每组基按顺序一一对应。由谱定理，这些基合起来就是全空间的标准正交基，因此得到的 $U$ 幺正；逐个基向量检查便有 $BU=UA$。所以在有限维自伴的条件下，分类数据是每个特征值及其重数。
\end{solution}

\section{忘却一部分结构与保留谱分量}

\begin{example}{忘掉算子后，允许的箭头变多了}
令 $F$ 将 $(H,A)$ 送到 $H$，将箭头 $U$ 送到同一个线性映射 $U$。证明 $F:\mathcal S\to\mathcal H$ 是函子。给出一个只在忘却后才允许的箭头。
\end{example}
\begin{solution}
一个\textbf{函子}要同时处理对象和箭头，并保持恒等箭头和复合。$\mathcal S$ 的每个箭头本来就是幺正同构，因此 $F(U)$ 的确是 $\mathcal H$ 的箭头；同一个恒等映射忘掉算子后仍是恒等映射。由于 $F$ 不改线性映射本身，$F(VU)=F(V)F(U)$。这就验证了函子条件。此处只忘掉附加的算子，故称为忘却函子。

取前例的 $(\mathbb C^2,D)$ 与 $(\mathbb C^2,A)$。两者经 $F$ 都变成 $\mathbb C^2$，而 $I$ 当然是 $\mathcal H$ 中的箭头。若这个 $I$ 来自原对象之间的箭头，就须有 $AI=ID$，即 $A=D$，矛盾。这说明在这对对象之间，忘却后的箭头多于原来箭头的像。术语上，$F$ \textbf{不满}；但原来两条不同的箭头仍是两个不同的线性映射，故 $F$ 对每对对象的箭头集合是单射，称为\textbf{忠实}。
\end{solution}

\begin{example}{谱投影与换坐标的两条路径}
仍取 $AQ=QD$。令 $P_6(A)$、$P_6(D)$ 分别为到 $6$ 特征子空间的正交投影。证明 $P_6(A)Q=QP_6(D)$，并用 $y=(1,2)^T$ 验算。
\end{example}
\begin{solution}
多项式 $p(t)=(t-1)/5$ 在 $6$ 处取 $1$，在 $1$ 处取 $0$。在特征基上计算可知 $P_6(A)=p(A)$、$P_6(D)=p(D)$，于是
\[
 P_6(A)=\frac15\begin{pmatrix}4&2\\2&1\end{pmatrix},\qquad
 P_6(D)=\begin{pmatrix}1&0\\0&0\end{pmatrix}.
\]
由 $AQ=QD$ 得 $(A-I)Q=Q(D-I)$，即所求等式。具体地，$Qy=(0,\sqrt5)^T$，所以 $P_6(A)Qy=(2,1)^T/\sqrt5$；另一条路径先得 $P_6(D)y=(1,0)^T$，再乘 $Q$ 也得 $(2,1)^T/\sqrt5$。

一般地，交织箭头把每个特征子空间幺正地送到同一特征值的子空间，并保持其正交补，故 $P_\lambda(B)U=UP_\lambda(A)$（没有该特征值时取零投影）。把每个对象 $(H,A)$ 忘却为向量空间 $H$，记这一函子为 $G$；它的目标范畴允许所有线性映射。上述交换条件恰好说 $P_\lambda$ 是从 $G$ 到 $G$ 的一个\textbf{自然变换}：每个空间上的分量是一个投影，且这些分量与所有允许的换坐标相容。投影往往不可逆，因此这里的目标范畴不能只允许幺正同构。
\end{solution}

\section{正交和与选择的不唯一性}

\begin{example}{正交直和如何合并重复特征值}
把 $A$ 与一维算子 $[6]$ 做正交直和。求 $A\oplus[6]$ 的特征子空间，并构造一个混合两块而仍与它交织的幺正变换。
\end{example}
\begin{solution}
正交直和的内积是两块内积之和。若两块各有交织箭头 $U,V$，则 $U\oplus V$ 保持这个内积，并且
\[
 (B\oplus B')(U\oplus V)=(U\oplus V)(A\oplus A')
\]
由两块各自的交织等式逐项得到。对本题，$Q\oplus[1]$ 给出对角形 $D\oplus[6]=\operatorname{diag}(6,1,6)$ 到 $A\oplus[6]$ 的箭头。设 $q_1,q_2$ 是 $Q$ 的两列，$v=(q_1,0)$、$w=(0,0,1)^T$、$s=(q_2,0)$。则 $6$ 的特征子空间由 $v,w$ 张成，$1$ 的由 $s$ 张成。特征值重数按块相加。

令 $Rv=(v+w)/\sqrt2$、$Rw=(-v+w)/\sqrt2$、$Rs=s$。三个像仍是标准正交基，所以 $R$ 幺正。算子 $A\oplus[6]$ 在 $v,w$ 所张成的空间上就是 $6I$，在 $s$ 上就是 $I$；分别在这三个基向量上比较两种作用顺序，可得 $R(A\oplus[6])=(A\oplus[6])R$。因此来自不同块的同一特征值分量可以互相混合。这里的“正交和”是对对象和箭头的逐块构造；通常的坐标嵌入并非幺正同构，不能直接把它当成这个群胚中的余积结构。
\end{solution}

\begin{example}{能否自然地选定一个单位特征向量}
假设要在每个非零的 $(H,A)$ 中选一个单位特征向量 $q_A$，并要求对任意交织箭头 $U:(H,A)\to(K,B)$，都有 $Uq_A=q_B$。这样的选择是否可能？
\end{example}
\begin{solution}
只看一维对象 $(\mathbb C,[6])$。它的每个单位复数都是特征向量。取 $U=-I$，则 $U$ 幺正且和 $[6]$ 交换，所以是此对象到自身的箭头。相容要求此时成为 $-q_A=q_A$，只能推出 $q_A=0$，和单位长度矛盾。这说明即使特征值只有一个、重数也是一，也不能按上述要求自然地挑选特征向量，更不能自然地指定一整组有序特征基。

另一方面，这个对象上的 $6$ 谱投影就是恒等映射，不因 $U=-I$ 改变。多维中也一样：把单位特征向量 $q$ 换成 $e^{i\theta}q$，投影 $qq^*$ 不变；有重数时，在同一子空间内更换标准正交基，也不改变到整个子空间的投影，因为投影由“在子空间上恒等、在正交补上为零”唯一确定。谱定理保证特征基\textbf{存在}，而真正与换坐标相容的是特征子空间及其谱投影。
\end{solution}

\appendix
\chapter{综合例题：把七个视角连起来}
\label{ch:synthesis}

做综合题时，先问对象是什么：平面点集、二次型、线性算子，还是带算子的内积空间。再问允许的箭头是什么。这样就能判断哪些量应该保持，以及哪些反推需要额外条件。

\section{从几何图形反推代数数据}

\begin{example}{先平移，再旋转，最后求面积}
求曲线
\[
5x^2+4xy+2y^2-6x-1=0
\]
的中心、半轴长度及围成区域的面积，并写出从标准椭圆到原椭圆的全等变换。
\end{example}
\begin{solution}
把左边记作 $f(X)$，$X=(x,y)^T$。其梯度为 $2A_0X-2b$，其中 $b=(3,0)^T$。令 $A_0h=b$，解得 $h=(1,-1)^T$，因为 $A_0(1,-1)^T=(3,0)^T$。写 $X=h+Z$，展开后一次项抵消：
\[
f(h+Z)=Z^TA_0Z-h^TA_0h-1=Z^TA_0Z-4.
\]
所以中心为 $h$。再令 $Z=Q_0(u,v)^T$，标准方程为 $6u^2+v^2=4$。半轴长为 $2/\sqrt6$ 和 $2$，围成面积为 $4\pi/\sqrt6$。全等对应是 $(u,v)^T\mapsto h+Q_0(u,v)^T$；平移与正交变换都保距离。Hessian $2A_0$ 说明这是正定二次型的等值线，实分析的变量代换则给出相同的面积结论。
\end{solution}

\begin{example}{变换前后谱不同，为什么图形仍完全相同}
对比三组方程：
\[
\begin{array}{ll}
\text{甲：}&x^2+4y^2=1,\quad 2x^2+8y^2=2;\\
\text{乙：}&x^2-4y^2=0,\quad 2x^2-8y^2=0;\\
\text{丙：}&x^2+4y^2=1,\quad 2x^2+8y^2=1.
\end{array}
\]
判断点集是否全等，并说明为何不与同谱判据冲突。
\end{example}
\begin{solution}
甲组两边同除以 $2$ 后方程相同。正定椭圆若写为 $x^TAx=c$ 且 $c>0$，应把它化为 $x^T(A/c)x=1$，用于判断的是归一化矩阵 $A/c$ 的谱。甲组归一化矩阵相同。乙组的两条方程都是 $x=\pm2y$，点集也相同；但零水平不能除以 $c$ 归一化，正定单位椭圆的判据不适用。丙组两个右端都为 $1$，归一化谱不同，半长轴依次为 $1$ 和 $1/\sqrt2$，故不全等。三组对照说明：谱信息、等值常数、二次型的类型都要放在同一个命题中核对。
\end{solution}

\section{从离散谱走向连续谱}

\begin{example}{一串有特征值的算子，极限却没有特征值}
在 $L^2[0,1]$ 中，令 $s_n(t)=k/n$ 当 $k/n\le t<(k+1)/n$（$k=0,\ldots,n-1$），端点 $1$ 任意赋值。令 $M_nf=s_nf$，$Mf=tf$。求 $M_n$ 的谱，证明 $M_n$ 在算子范数下趋于 $M$，并解释极限为何可没有特征值。
\end{example}
\begin{solution}
将 $[0,1]$ 分成 $n$ 段，每段上都有非零的平方可积函数。支撑在第 $k$ 段的非零函数都是 $M_n$ 的特征函数，特征值为 $k/n$。若 $\lambda$ 不属于这 $n$ 个数，乘以 $1/(s_n-\lambda)$ 是 $(M_n-\lambda I)$ 的有界逆，因此谱恰为这 $n$ 个值。

由 $|s_n(t)-t|\le1/n$ 几乎处处，有 $\|(M_n-M)f\|_2\le \|f\|_2/n$。所以 $M_n\to M$。有限个谱点越来越密，在极限填满 $[0,1]$；对极限算子，要直接检查有界可逆性：$[0,1]$ 外可乘 $1/(t-\lambda)$ 取逆，区间内可用集中在 $\lambda$ 附近的小区间上的单位函数使 $\|(M-\lambda I)f\|_2$ 任意小。可是 $Mf=\lambda f$ 要求 $f$ 支撑在 $\{t=\lambda\}$ 上，这个集合测度为零，只能得到零函数。每个近似算子的特征空间支撑在一整段上；极限中单个取值点不再承载非零函数。
\end{solution}

\Needspace{9\baselineskip}
\section{复空间与实空间之间}

\begin{example}{一个复特征方向变成两个实方向}
令 $H=\begin{pmatrix}2&i\\-i&2\end{pmatrix}$ 作用在 $\mathbb C^2$ 上。求酉标准形，并把 $z^*Hz=1$ 看作 $\mathbb R^4$ 的全等问题。
\end{example}
\begin{solution}
$H^*=H$。分别计算可得
\[
H\binom{i}{1}=3\binom{i}{1},\qquad
H\binom{-i}{1}=\binom{-i}{1}.
\]
这两个向量在复内积下正交，各有长度 $\sqrt2$。把归一化向量组成列矩阵 $U$，有 $U^*U=I$、$U^*HU=\operatorname{diag}(3,1)$。令 $z=Uw$、$w_j=a_j+ib_j$，则
\[
z^*Hz=3(a_1^2+b_1^2)+(a_2^2+b_2^2).
\]
把每个复坐标拆成实部、虚部，复范数平方恰是四个实坐标的平方和，故 $U$ 同时给出 $\mathbb R^4$ 的一个正交变换。于是原等值集与上述实四维椭球全等，半轴长度为 $1/\sqrt3,1/\sqrt3,1,1$。一个一维复特征空间有两个实维度，这就是半轴成对出现的原因。

这里还有一层结构：复线性变换与乘 $i$ 交换；实正交变换未必如此。例如 $z\mapsto\overline z$ 在实坐标上保距，却把 $iz$ 送成 $-i\overline z$。忘掉复结构会允许更多等距箭头，这与范畴论章的忘却操作相呼应。
\end{solution}

\end{document}
